% ============================================================================
% CopilotScope -- the manual: what it is, how it measures, and how to use it.
%
% This is NOT a research paper. research/articles/quality_measurement_framework.tex
% is the formal treatment of the scoring model; this document is the manual a
% person reads to understand the system and then operate it, and it carries a
% tutorial for exactly that reason.
%
% Sources of truth, which this file presents rather than extends: README.md,
% docs/TUTORIAL.md, docs/SIGNAL_COVERAGE.md, docs/PRIVACY.md,
% docs/CALIBRATION.md, SECURITY.md and GOVERNANCE.md. It introduces no fact
% those documents do not already contain; where it disagrees with them, they
% are right and this is stale.
%
% The diagrams are not redrawn here. They are rendered from docs/diagrams/*.mmd
% -- the same files docs/DIAGRAMS.md embeds, kept identical by
% scripts/check-diagrams.mjs -- into vector PDFs by scripts/render-diagrams.mjs.
% One source, two output formats, and both language editions share the result.
%
% A Polish edition lives alongside this file:
%   copilotscope-manual.pl.tex
% scripts/check-doc-parity.mjs fails the build if one is edited without the
% other.
%
% Build (PDFs are build output -- never committed):
%   npm ci && npm run render:diagrams
%   pdflatex copilotscope-manual.tex && pdflatex copilotscope-manual.tex
% Also built by .github/workflows/build-docs-pdf.yml, which builds both editions.
% ============================================================================
\documentclass[11pt,a4paper]{article}

\usepackage[utf8]{inputenc}
\usepackage[T1]{fontenc}
\usepackage{lmodern}
\usepackage{geometry}
\usepackage{xcolor}
\usepackage{hyperref}
\usepackage{graphicx}
\usepackage{enumitem}
\usepackage{fancyhdr}
\usepackage{titlesec}
\usepackage{parskip}
\usepackage{booktabs}
\usepackage{array}
\usepackage{tabularx}
\usepackage{longtable}
\usepackage{fancyvrb}
\usepackage{microtype}
\usepackage{amssymb}
\usepackage{url}
\urlstyle{tt}

\geometry{a4paper, left=2.2cm, right=2.2cm, top=2.5cm, bottom=2.5cm}

% House colors
\definecolor{primaryblue}{RGB}{41,128,185}
\definecolor{darkblue}{RGB}{31,78,121}
\definecolor{lightgray}{RGB}{236,240,241}
\definecolor{darkgray}{RGB}{52,73,94}
\definecolor{accentgreen}{RGB}{39,174,96}
\definecolor{accentorange}{RGB}{230,126,34}
\definecolor{accentred}{RGB}{231,76,60}

\titleformat{\section}
    {\color{primaryblue}\Large\bfseries\sffamily}
    {\thesection.}{0.5em}{}[\titlerule]
\titleformat{\subsection}
    {\color{darkblue}\large\bfseries\sffamily}
    {\thesubsection}{0.5em}{}
\titleformat{\subsubsection}
    {\color{darkgray}\normalsize\bfseries\sffamily}
    {\thesubsubsection}{0.5em}{}

\newcommand{\docstatus}{NOT CALIBRATED --- READ SECTION 7}

\pagestyle{fancy}
\fancyhf{}
\renewcommand{\headrulewidth}{0.4pt}
\fancyhead[L]{\small\sffamily\color{darkgray} CopilotScope --- the manual}
\fancyhead[R]{\small\sffamily\color{accentred} \docstatus}
\fancyfoot[C]{\thepage}

\hypersetup{
    colorlinks=true,
    linkcolor=primaryblue,
    urlcolor=primaryblue,
    pdfauthor={Konrad Cinkusz},
    pdftitle={CopilotScope: quality scoring for AI coding-assistant sessions}
}

% Inline code. \detokenize makes every character literal, so an underscore in a
% configuration key or a backslash in a Windows path needs no escaping at the
% call site -- which is the single most common way a document like this fails
% to build.
\newcommand{\code}[1]{\texttt{\detokenize{#1}}}
% A path or a long identifier inside a narrow table column. \path breaks after
% / . _ and - , which \texttt does not -- long project paths otherwise run past
% the column and produce an overfull box per row.
\newcommand{\codepath}[1]{{\small\path{#1}}}
% Same breaking behaviour at body size, for a long identifier sitting in prose
% where \code would push the line past the margin.
\newcommand{\codewrap}[1]{\path{#1}}
\newcommand{\finding}[1]{{\color{primaryblue}\textbf{#1}}}

\DefineVerbatimEnvironment{shell}{Verbatim}
    {frame=single, rulecolor=\color{darkgray}, framesep=5pt,
     fontsize=\small, xleftmargin=2pt, xrightmargin=2pt}

\begin{document}

\thispagestyle{plain}
\begin{center}
    {\LARGE\bfseries\sffamily\color{darkblue} CopilotScope}

    \vspace{0.3cm}

    {\Large\sffamily\color{darkgray} Quality scoring for AI coding-assistant
    sessions --- what it measures, how it works, and how to run it}

    \vspace{0.6cm}

    {\large Konrad Cinkusz}

    \vspace{0.2cm}

    {\normalsize\color{darkgray} \url{https://github.com/konradcinkusz/copilot-scope}}

    \vspace{0.2cm}

    {\normalsize\color{accentred}\sffamily \docstatus\ --- \today\ ---
     \textit{Polska wersja: } \texttt{copilotscope-manual.pl.tex}}
\end{center}

\vspace{0.4cm}

\begin{center}
\begin{minipage}{0.92\linewidth}
\small
\textbf{What this document is.} Everything needed to understand CopilotScope and
then operate it: the problem it addresses, how it is built, exactly what it
measures and what it refuses to measure, and a tutorial that takes a reader from
an empty machine to a scored session, a team deployment and a Grafana panel.

\textbf{What it is not.} It is not a research paper --- that is
\texttt{research/articles/quality\_measurement\_framework.tex} --- and it is not
a validation of the score. The composite has not been calibrated against human
labels. Section~7 says so in full, because a number presented without its
limits is worse than no number.
\end{minipage}
\end{center}

\vspace{0.5cm}

{\small\tableofcontents}

\newpage

% ============================================================================
\section{The problem, and what this measures instead}
% ============================================================================

Every vendor of an AI coding assistant will tell you how much of it your team
used. Seats, suggestions shown, suggestions accepted, tokens burned. Those
numbers are real, easy to collect, and answer a question nobody actually has.
The question a team has is whether the tooling is helping --- and usage cannot
answer it, because usage goes up in exactly the same way when the assistant is
working well and when a developer is stuck in a loop asking the same thing four
different ways.

\finding{CopilotScope scores sessions, not usage.} A session is one conversation
with an assistant. For each one it produces a composite quality score with a
confidence figure beside it, the specific turn where the conversation stalled
and why, whether the edits the developer accepted actually survived in the file,
and what the repair loop cost in tokens and wall-clock time.

Three properties shape everything that follows:

\begin{enumerate}[leftmargin=*]
    \item \textbf{Nothing leaves the machine.} The collector runs on your
    infrastructure and has no callback of any kind. This is not a policy
    statement; there is no code path that phones anywhere.
    \item \textbf{The formula is published and deterministic.} Six weighted
    components, renormalized over the ones that reported data, with confidence
    exported next to every score. You can re-derive any number by hand.
    \item \textbf{It cannot rank developers.} There is no per-developer
    dimension in any view or export, and tests assert that there is not. This is
    enforced rather than promised, and Section~7.3 explains why that
    constraint is load-bearing rather than decorative.
\end{enumerate}

\subsection{Who this is for}

A team lead asking whether last month's model upgrade helped or hurt. A platform
engineer who already runs Prometheus and wants quality signals in it. A
developer who wants to know which of their sessions went badly and at which
turn. An organisation that needs the measurement to be defensible to a works
council before it is allowed to exist at all.

It is not for someone who wants a leaderboard. That is not a limitation to be
worked around; see Section~7.3.

\subsection{What a session looks like once it has been measured}

\begin{figure}[htbp]
\centering
\includegraphics[width=0.62\linewidth]{../diagrams/rendered/a1-system-context.pdf}
\caption{The system boundary. Four assistants emit OpenTelemetry to one endpoint;
the collector tells them apart from the payload rather than from configuration.
The GitHub Metrics API is the exception --- it is polled, and what it returns is
org-level usage, context beside the score rather than the score itself.}
\label{fig:context}
\end{figure}

% ============================================================================
\section{Architecture}
% ============================================================================

\subsection{What each project is}

\begin{figure}[htbp]
\centering
\includegraphics[width=\linewidth]{../diagrams/rendered/a2-solution-layout.pdf}
\caption{The collector is the product. Everything else is a view onto it, a
development convenience, or opt-in.}
\label{fig:layout}
\end{figure}

\begin{longtable}{@{}p{0.30\linewidth}p{0.44\linewidth}p{0.18\linewidth}@{}}
\toprule
\textbf{Project} & \textbf{Role} & \textbf{Dependencies} \\
\midrule
\endhead
\codepath{src/CopilotScope.Collector} & OTLP/HTTP ingest with an in-repo protobuf
decoder, session aggregation, the quality engine, turn analysis, the insight
pipeline, the REST API, the Prometheus exporter, persistence & Npgsql only \\
\codepath{src/CopilotScope.Dashboard} & Blazor Server UI: sessions, quality meter,
turn analysis, transcript, team views & none \\
\codepath{src/CopilotScope.ServiceDefaults} & Shared kernel: self-instrumentation,
health, service discovery, HTTP resilience. No domain logic & OpenTelemetry,
Microsoft.Extensions \\
\codepath{src/CopilotScope.AppHost} & Aspire orchestration for development:
Postgres, pgAdmin, both services & Aspire.Hosting \\
\codepath{src/CopilotScope.JudgeAgent} & Opt-in LLM judge: G-Eval, SPUR, RAGAS,
deep friction, task completion. Azure AI Foundry or any OpenAI-compatible
endpoint & Azure.AI, Microsoft.Agents.AI \\
\codepath{src/CopilotScope.AgentForge} & Opt-in persona agents grounded on consented
transcripts & Azure.AI, Microsoft.Agents.AI \\
\codepath{tools/CopilotScope.Seeder} & Fabricated demo sessions pushed into a
running collector & none \\
\codepath{tools/CopilotScope.LogImporter} & Claude Code transcripts parsed into
scored sessions & none \\
\codepath{tools/CopilotScope.TelemetryGen} & One realistic session over real
OTLP/HTTP, used as a probe & none \\
\bottomrule
\caption{The solution, project by project. The three tools ship together in one
\code{copilotscope-tools} container image, so none of them requires a .NET SDK
to run.}
\end{longtable}

The collector pulls in \emph{one} NuGet dependency at runtime, and the dashboard
none at all. That is deliberate: the OTLP protobuf decoder is written in-repo
against the official \code{.proto} field numbers rather than taking the
OpenTelemetry SDK, and the dashboard polls the collector's REST API every two
seconds rather than carrying a JavaScript framework.

\subsection{What happens to one batch of telemetry}

\begin{figure}[htbp]
\centering
\includegraphics[width=0.72\linewidth]{../diagrams/rendered/a3-ingest-pipeline.pdf}
\caption{Every arrow that leaves the happy path logs why it left. This diagram is
the answer to ``the stack is running and the dashboard is empty''.}
\label{fig:ingest}
\end{figure}

Two decisions in that pipeline are worth stating explicitly, because both are
the kind that only look important after they have gone wrong somewhere.

\textbf{Redaction happens before aggregation.} Under privacy mode, identifying
values are removed from the batch before anything downstream --- the store, the
write-behind snapshot, the Prometheus exporter --- has seen them. The guarantee
therefore survives a database dump, rather than depending on every read path
remembering to apply it.

\textbf{The decoded payload is bounded.} Kestrel limits the compressed request
size; gzip reaches roughly 1000:1. An unbounded copy of a compressed body is a
memory-exhaustion vector, and on a collector with no key configured it is
reachable before authentication. The ceiling is 64~MB per batch.

\subsection{Where the data lives, and what eviction does not mean}

Sessions live in two places at once. Ingest marks a session dirty; a background
writer upserts the dirty set once per second, so a burst of telemetry does not
become a burst of writes. The full session state goes to Postgres as a
\code{jsonb} snapshot --- counters, time-to-first-token samples, tool statistics,
per-turn aggregates, edits, feedback, the event tail and the captured transcript
--- with a few columns pulled out for querying.

\begin{figure}[htbp]
\centering
\includegraphics[width=0.78\linewidth]{../diagrams/rendered/a4-session-lifecycle.pdf}
\caption{The in-memory store is a bounded working set, not the extent of your
history. Eviction is not deletion.}
\label{fig:lifecycle}
\end{figure}

Reads are served from Postgres with the live in-memory aggregates layered on
top, so the API answers over the whole archive while a session being typed into
right now is still current. Without a connection string the collector degrades
to memory alone and reports \code{durable: false} rather than pretending
otherwise --- which is also why \code{dotnet run} on a bare machine works with
no database at all.

% ============================================================================
\section{Getting data in}
% ============================================================================

\subsection{The install, end to end}

\begin{figure}[htbp]
\centering
\includegraphics[width=0.70\linewidth]{../diagrams/rendered/b1-install-and-connect.pdf}
\caption{Two commands, plus one manual step per assistant that no script can
take for you.}
\label{fig:install}
\end{figure}

On one machine there is nothing to declare: no key to generate, no environment
variable to export, no JSON to hand-edit. Every published port binds to
\code{127.0.0.1}, Postgres publishes no port at all and trusts its unpublished
socket, and the ingest key is empty, which is the collector's open mode. A
credential on loopback protects nothing and costs a configuration step in every
client --- and that step is the one people get wrong.

\subsection{Why configuration used to be the hard part}

\begin{figure}[htbp]
\centering
\includegraphics[width=\linewidth]{../diagrams/rendered/b2-where-each-switch-lives.pdf}
\caption{Four surfaces, three kinds of place to configure them, and only one of
those places survives opening a new terminal.}
\label{fig:switches}
\end{figure}

Claude Code is the best case. An \code{env} block in \code{~/.claude/settings.json}
applies to every terminal, every project, and to Claude Code running inside VS
Code. Four values matter, and the first two are the ones people omit:

\begin{itemize}[leftmargin=*]
    \item \code{CLAUDE_CODE_ENABLE_TELEMETRY=1} --- the master switch. Without
    it nothing at all is exported, whatever else is set.
    \item \code{OTEL_LOGS_EXPORTER=otlp} --- the log events are what carry the
    session. A default install emits metrics and events and no spans, so a
    configuration that sets only the metrics exporter produces an almost empty
    session.
    \item \code{OTEL_EXPORTER_OTLP_PROTOCOL=http/protobuf} and the endpoint.
\end{itemize}

Content capture here is three separate opt-ins
(\codewrap{OTEL_LOG_USER_PROMPTS}, \codewrap{OTEL_LOG_ASSISTANT_RESPONSES},
\codewrap{OTEL_LOG_TOOL_DETAILS}), and the OpenTelemetry GenAI standard variable
that works for Copilot CLI is \emph{not} read by Claude Code. Time to first
token needs the tracing beta, which is the only source of that signal for this
assistant.

VS Code is configured through its user settings and needs a window reload,
because the extension reads its configuration at startup. Copilot CLI reads
environment variables and nothing else. Cowork is configured in the desktop
app's own settings UI, wants the full \code{/v1/logs} path rather than the base
endpoint, requires a Team or Enterprise plan and org admin access, and exports
log events only.

\subsection{The path that needs no configuration at all}

\begin{figure}[htbp]
\centering
\includegraphics[width=0.72\linewidth]{../diagrams/rendered/b3-transcript-import.pdf}
\caption{Claude Code writes every session to disk whether or not telemetry is
configured --- and most developers never configure it.}
\label{fig:import}
\end{figure}

Two refusals in that path are deliberate and worth defending. A repository label
is left \emph{absent} rather than guessed from a directory name, because a guess
would invent a second cohort for a project the collector already knows by its
git remote --- two rows for one repository, which is worse than no label. And a
session already held from live telemetry is not overwritten by an import,
because the transcript carries less signal and would quietly lower a score that
was measured on more evidence.

Imported sessions are badged and carry genuinely lower confidence. The
transcript records tokens, models, tool calls and outcomes, turn boundaries and
real wall-clock timings. It does not record time to first token, edit
accept/reject decisions or thumbs feedback, because those are OpenTelemetry
events. Rather than defaulting them to zero --- which would read as ``every edit
was rejected'' --- they are left absent, and the composite renormalizes.

% ============================================================================
\section{The measurement}
% ============================================================================

\subsection{The composite score}

\begin{figure}[htbp]
\centering
\includegraphics[width=0.80\linewidth]{../diagrams/rendered/c1-composite-score.pdf}
\caption{Six components. Only the ones that reported data enter the composite,
and the weights are renormalized across them.}
\label{fig:composite}
\end{figure}

\begin{center}
\begin{tabular}{@{}llp{0.52\linewidth}@{}}
\toprule
\textbf{Component} & \textbf{Weight} & \textbf{What it is} \\
\midrule
Reliability & 0.25 & Squared error-free rate across LLM and tool calls. Squared,
so a session with a few errors is penalised more than linearly \\
Acceptance  & 0.20 & Accepted edits, paired with edit survival so that accepting
bad suggestions does not read as success \\
Friction    & 0.20 & Mean turn score from the turn analyser \\
Latency     & 0.15 & How much of the session sat past the 2~s attention and
8~s abandonment thresholds \\
Feedback    & 0.10 & Thumbs, where the assistant reports them \\
Efficiency  & 0.10 & Token economics: cost per turn and per accepted edit \\
\bottomrule
\end{tabular}
\end{center}

\vspace{0.3cm}

\finding{Renormalization is the part to understand.} Only components with actual
data enter the composite, and the weights are rescaled across them. An earlier
version let empty components contribute a neutral prior, with the result that
almost every session scored near 80 --- the measurement was reporting the prior,
not the session. Confidence, exported next to every score, is data coverage
multiplied by a sample ramp: a 90 built on four samples means less than a 70
built on forty.

Grades are thresholds on the composite: 85 and above excellent, 70 good, 55
fair, 40 poor, below that critical. A percentile rank is computed against a
fixed trailing window of user sessions rather than against whatever happened to
be in memory, which is what makes it comparable over time.

\subsection{Turn analysis, or: which turn went wrong}

\begin{figure}[htbp]
\centering
\includegraphics[width=0.88\linewidth]{../diagrams/rendered/c2-turn-analysis.pdf}
\caption{A composite score says a session was poor. This says which turn, and
why.}
\label{fig:turns}
\end{figure}

Every \code{invoke_agent} trace is one turn. Each turn is scored for LLM and
tool errors, for latency against \emph{this session's own} median time to first
token, and for repair loops --- bursts of tool calls containing failures. The
dashboard surfaces the best and worst turn with the reasons attached.

The per-session median is the load-bearing choice. A global latency threshold
cannot tell a slow turn from a slow model, or a slow model from a slow morning;
a session compared against itself can.

Note the consequence for assistants that do not emit a turn boundary: Cursor's
export sends metrics and logs but no traces, and a turn here \emph{is} a trace,
so turn-level friction analysis cannot run for it at all even if the export
works.

\subsection{The insight pipeline}

Analyzers are pluggable: one \code{IInsightAnalyzer} class plus one dependency
injection registration adds an algorithm with no UI work. The in-process ones
are edit survival, throughput, latency utility, token economics and workflow
friction.

The LLM-graded algorithms are deliberately \emph{not} in that pipeline. The
interface is synchronous and in-process; a judge call is an asynchronous network
call to a model server, with a cost and a rate limit. They live in the separate,
opt-in judge service instead, reachable at
\code{POST /api/sessions/{id}/judge}, and that service can be pointed at Azure
AI Foundry or at an OpenAI-compatible endpoint on your own hardware --- so a
self-hosted deployment can run them without transcripts leaving the building.

% ============================================================================
\section{The ten algorithms, and what each one is worth today}
% ============================================================================

The project tracks ten evaluation algorithms on two axes: whether the
implementation exists, and whether it runs locally or needs model access.

\begin{longtable}{@{}p{0.03\linewidth}p{0.26\linewidth}p{0.20\linewidth}p{0.42\linewidth}@{}}
\toprule
\textbf{\#} & \textbf{Algorithm} & \textbf{Status} & \textbf{Notes} \\
\midrule
\endhead
1 & LLM-as-a-Judge (G-Eval) & implemented, opt-in & Needs an LLM endpoint; local
model works \\
2 & SPUR satisfaction rubrics & implemented, zero-shot & Directional until
labelled sessions exist to calibrate against \\
3 & RAG component metrics (RAGAS) & implemented, opt-in & Only meaningful for
retrieval flows; reports no-data until retrieval context is captured \\
4 & Edit survival & full, local & Four-gram and no-revert split, 0.4/0.6; also
feeds the acceptance component \\
5 & Acceptance-weighted throughput & full, local & Accepted lines per turn and
per 1k output tokens, rejection-discounted \\
6 & Turn-level friction and repair & full, local & The turn analyser; also the
friction component \\
7 & Latency-utility model & full, local & Per-sample utility curve with the 2~s
and 8~s risk buckets \\
8 & Token and cache economics & full, local & Per-model cost, cache savings,
cost per turn and per accepted edit \\
9 & Workflow-friction signals & simplified local, deep in the judge & Counts
observed repair events, not inferred emotion. Off by default, report-only,
never scored \\
10 & Task-completion detection & partial local, implemented in the judge & No
local ingest path for build and test exit codes yet \\
\bottomrule
\caption{Implementation status. ``Opt-in'' means configuration, not whether the
container runs.}
\end{longtable}

Algorithm 9 deserves its own sentence, because it is the one most likely to be
misread. It counts \emph{observed repair events} --- rephrasing measured by
Jaccard similarity between consecutive prompts, corrective replies, negative
feedback markers --- and not emotion. It is excluded from the composite on
purpose: a lexical heuristic does not belong in a number somebody might act on.
It is also the only analyzer that reads the developer's own prompt text, which
is why it stays switched off until an operator turns it on, and why quoting
messages needs a second, separate opt-in.

% ============================================================================
\section{Signal coverage: why an 80 here is not an 80 there}
% ============================================================================

\begin{figure}[htbp]
\centering
\includegraphics[width=0.92\linewidth]{../diagrams/rendered/c3-signal-coverage.pdf}
\caption{Every assistant lands in the same schema. They do not report the same
signals.}
\label{fig:coverage}
\end{figure}

All four supported surfaces normalize onto one internal vocabulary, so
aggregation, scoring and the dashboard treat them identically. What differs is
what each one actually sends.

\begin{center}
\begin{tabular}{@{}lcccc@{}}
\toprule
\textbf{Signal} & \textbf{VS Code} & \textbf{Copilot CLI} & \textbf{Claude Code}
& \textbf{Cowork} \\
\midrule
Tokens, calls, errors      & yes & yes & yes & yes \\
Time to first token        & yes & yes & beta flag only & no \\
Edit accept / reject       & yes & no  & yes & yes \\
Thumbs feedback            & yes & no  & no  & no \\
Lines of code              & yes & no  & yes & no \\
Turn boundary (a trace)    & yes & yes & beta flag only & no \\
\bottomrule
\end{tabular}
\end{center}

\vspace{0.3cm}

\finding{Scores are comparable within an assistant and directional across them.}
A Claude Code session has no thumbs and, without the tracing beta, no
time-to-first-token, so its 80 rests on less evidence than a VS Code session's
80. The composite handles this honestly by renormalizing, and confidence reports
what it had to work with. What it cannot do is make the two numbers mean the
same thing --- so a bake-off between assistants has to hold the signal set
constant, not just the score.

Cursor is a special case and is documented as one. Its OpenTelemetry export is
Enterprise-plan only and sends metrics and logs but no traces. The collector
detects and normalizes \code{cursor.*} attributes, but no payload from a real
Cursor session has ever been tested against it, so it is listed as detected and
unverified rather than supported. That distinction exists because this project
previously claimed support for an assistant it had no captured payloads for.

% ============================================================================
\section{Honest limits}
% ============================================================================

This section exists because a measurement presented without its limits is worse
than no measurement: the reader trusts it.

\subsection{The composite is not calibrated}

There are no human labels in this repository and no calibration run has been
performed. The weights are a considered judgement, not a fitted model. Every
judge score is therefore directional and gates nothing.

The machinery to fix this exists and is waiting for input. The judge service
carries a calibration module that computes Cohen's $\kappa$ --- unweighted,
linear- and quadratic-weighted, each with a seeded bootstrap confidence interval
--- between the judge and a human panel, plus the panel's agreement with itself
as the ceiling no algorithm can climb past. Verdicts are \emph{calibrated},
\emph{not-calibrated}, \emph{ceiling-too-low} or \emph{insufficient-data}
against a configurable threshold, default $\kappa \geq 0.70$.

What was missing until recently was anything to consume. A labelling flow now
exists behind \codewrap{CopilotScope:Labelling:Enabled}: a rater panel on the
session view with a handle, a band from 0 to 3 per rubric with anchor text, and
a skip option that is stored as a real answer but not exported as a label.
\code{GET /api/labels/export} produces exactly the JSON the calibration engine
reads, pinned by a round-trip test.

Seeded sessions are excluded from export by default, and that exclusion is the
point: validating the scoring model against the synthetic personas this
repository wrote would be a circle, and a $\kappa$ built on one is worse than no
$\kappa$ at all.

\subsection{A single number is not a verdict}

Confidence is exported beside every score for a reason. Read the components, not
the headline. A session scoring 62 with high confidence and a clear worst turn
is more actionable than one scoring 88 on two samples.

\subsection{Not a developer scoreboard --- and this is enforced}

The score grades a \emph{session}, not a person. Ranking developers by it is a
misuse the design does not support: there is no per-developer dimension in the
cohort filter, no per-developer view, and no export with such a column. Tests
assert this.

On a team deployment the convention becomes machinery. Privacy mode
pseudonymizes identities at ingest, drops prompt and response content, refuses
to render any view covering fewer than $k$ developers, and logs every read.

The reasoning is not only ethical. A score used for appraisal stops measuring
the tooling almost immediately, because the people being measured adapt --- and
what they adapt is the thing being counted, not the thing you wanted. That is
also exactly why acceptance rate is only 0.20 of the composite and is paired
with edit survival: push on acceptance alone and you reward accepting bad
suggestions. Goodhart's law applies to this repository as much as to anything
it measures.

\subsection{What the project deliberately does not do}

It does not archive raw telemetry --- content capture stores a bounded preview,
the last hundred entries truncated at four thousand characters, and a full
archive is a job for a real telemetry backend, which the forwarder exists to
feed. It does not fold delivered outcomes into the score, even though it now
links sessions to the pull requests they produced, because the correlation study
that would justify doing so is the reason the outcomes are being collected. And
it does not promote the friction analyzer into the composite, for the reason
given in Section~5.

% ============================================================================
\section{Privacy, security and team deployments}
% ============================================================================

\subsection{The moment a second developer's telemetry arrives}

The moment a second person's sessions reach the same collector, you are running
a technical system capable of monitoring employee performance. In the EU that
triggers works-council co-determination on the capability alone, whatever you
intend to do with it. Privacy mode is what makes the intent enforceable rather
than stated.

\begin{itemize}[leftmargin=*]
    \item Identities are pseudonymized before anything stores them, with a salt
    you supply. An ephemeral salt is generated if you forget, and the collector
    warns loudly, because pseudonyms that change on every restart stop
    correlating history across a deploy.
    \item Prompt and response content is dropped at ingest regardless of how the
    clients are configured.
    \item No view renders for fewer than $k$ developers.
    \item Every read is logged to a table that the session retention sweep never
    touches, because the audit record must outlive the sessions it describes.
    \item Raw OTLP forwarding is refused unless the operator explicitly states
    that the upstream backend is covered by the same agreement --- a faithful
    relay of un-redacted payloads is a hole straight through every other
    control.
\end{itemize}

\code{GET /api/privacy} reports what a running deployment actually enforces, so
the answer to ``is it on?'' is not a matter of trusting a configuration file.

\subsection{The credential model}

\begin{figure}[htbp]
\centering
\includegraphics[width=0.80\linewidth]{../diagrams/rendered/d1-deployment-modes.pdf}
\caption{The zero-credential default is scoped, not lax. Publishing beyond
loopback is a different deployment and is treated as one.}
\label{fig:deploy}
\end{figure}

Gating follows the key's presence, in any environment: with
\codewrap{CopilotScope:Ingest:ApiKey} empty the collector runs open, and with a key
set, ingest, the query API and the Prometheus endpoint all require it.

For anything shared, split the one key into scopes. The key every developer's
editor holds should not also be the key that reads transcripts and deletes
history:

\begin{itemize}[leftmargin=*]
    \item \textbf{Ingest} --- \code{POST /v1/*} and nothing else. Held by every
    emitter.
    \item \textbf{Read} --- the query API and \code{/metrics}, which includes
    captured transcripts. Held by the dashboard and Prometheus.
    \item \textbf{Admin} --- destructive and administrative operations. Implies
    Read.
\end{itemize}

The dashboard gets its own sign-in with two roles, because transcripts are the
sensitive payload: a viewer sees scores, turns and aggregates; an admin also
sees conversation content and can delete. Neither compose file terminates TLS,
so anything shared needs a reverse proxy in front of it --- these credentials
travel in a header and a cookie.

The collector cannot work out for itself whether it is exposed: behind Docker
every request arrives from the bridge gateway, so the remote address proves
nothing. The compose files therefore pass the address they published on, and a
startup warning fires when that is not loopback and no key is set. Two commands
refuse the combination outright rather than warning about it.

% ============================================================================
\section{Tutorial}
% ============================================================================

Everything below is runnable. The only prerequisite is Docker; a .NET SDK is
needed exclusively for the contributor path in Section~9.9.

\subsection{Five minutes: from an empty machine to a scored session}

\subsubsection*{Step 1 --- install}

\begin{shell}
curl -fsSL https://raw.githubusercontent.com/konradcinkusz/\
copilot-scope/master/install.sh | sh
\end{shell}

On Windows, in PowerShell:

\begin{shell}
irm https://raw.githubusercontent.com/konradcinkusz/copilot-scope/\
master/install.ps1 | iex
\end{shell}

The installer checks Docker, writes the compose file and a \code{copilotscope}
control script into \code{~/.copilotscope}, starts the stack, waits for the
collector to answer, and offers to configure each assistant it finds. Pass
\code{--yes} after \code{| sh -s --} to accept every offer without being asked.

Nothing else is required. There is no key to generate and no \code{.env} to
fill in.

\subsubsection*{Step 2 --- prove the pipeline works before blaming a client}

\begin{shell}
copilotscope probe
\end{shell}

This sends one simulated session over real OTLP/HTTP protobuf and then checks
that the collector can serve it back. If it passes, ingest, decoding, scoring
and persistence all work, and anything still missing is client configuration.
Doing this first saves an hour of debugging the wrong half of the system.

\subsubsection*{Step 3 --- look at something}

\begin{shell}
copilotscope demo          # fabricated sessions, badged demo
copilotscope open          # http://localhost:5200
\end{shell}

Seeded sessions are clearly marked, so a screenshot of them is never mistaken
for evidence about a real assistant. \code{copilotscope demo demo} loads the
larger multi-day dataset instead of the quick one.

\subsection{Connecting each assistant}

\begin{shell}
copilotscope connect claude-code     # ~/.claude/settings.json
copilotscope connect vscode          # VS Code user settings
copilotscope connect copilot-cli     # shell rc, or Windows User scope
copilotscope connect cowork          # prints what to type into the app
copilotscope connect all             # Claude Code and VS Code together
\end{shell}

Two flags apply to all of them:

\begin{itemize}[leftmargin=*]
    \item \code{--capture} also exports prompt and response text. Off by
    default, and it is the switch that turns a metadata-only deployment into one
    holding source code and anything a developer pasted into a chat.
    \item \code{--print} shows exactly what would be written and changes
    nothing.
\end{itemize}

\code{--traces} is Claude Code only and turns on the tracing beta, which is the
only source of time to first token for that assistant. The span schema can still
change; that is what the beta flag means.

Then the manual step the tooling cannot take for you: reload the VS Code window,
restart Claude Desktop for Cowork, or open a new terminal for Copilot CLI.
Finally, chat in \emph{agent mode} --- inline completions alone produce no chat
telemetry, which is the second most common reason for an empty dashboard.

\code{copilotscope disconnect <target>} removes exactly the keys that were
added, leaving the rest of your settings untouched.

\subsection{Scoring the history you already have}

If you use Claude Code, skip all of the above and run:

\begin{shell}
copilotscope import --dry-run     # see what it found, send nothing
copilotscope import               # import it
\end{shell}

Re-running is safe: sessions keep Claude Code's own identifier, so a second run
replaces rather than duplicates. Put it on a schedule if you like. Prompt text
stays out unless \code{--include-content} is passed.

One caveat the command states itself: run from the container, the importer
cannot reach git, so imported sessions carry no repository label. When that
matters, run it from a clone with a .NET SDK instead.

\subsection{Reading a session}

Open a session and read it in this order.

\begin{enumerate}[leftmargin=*]
    \item \textbf{Grade and confidence together.} Never the grade alone. Low
    confidence means the session did not emit enough for the number to carry
    weight.
    \item \textbf{The worst turn, with its reasons.} This is the actionable
    part: an error, a latency outlier against this session's own median, or a
    repair loop.
    \item \textbf{The components.} A composite of 62 driven by latency is a
    different problem from a 62 driven by reliability, and they have different
    fixes.
    \item \textbf{Edit survival, if the assistant reports it.} Acceptance
    without survival is the signature of code accepted and then reverted.
    \item \textbf{Token economics.} Cost per accepted edit is the number that
    turns a quality conversation into a budget conversation.
\end{enumerate}

The built-in documentation page at \code{/docs} explains every tile and every
component in interpretation language rather than metric language.

\subsection{Comparing windows and cohorts}

Two questions get their own endpoints, both filterable by repository, assistant
and model, and both exportable as CSV with group rows only:

\begin{shell}
GET /api/cohorts?days=30      # roll the window up by each axis
GET /api/compare?days=30      # this window against the one before it
GET /api/digest?days=7        # the week, as something to forward
\end{shell}

\code{/api/compare} is the ``did the upgrade help'' question, with the baseline
defaulting to the equally long window immediately before the current one.

Every axis describes the tooling. None of them describes a developer, and that
is not an oversight.

\subsection{Prometheus and Grafana}

If you already run them, the collector exposes its \emph{derived} signals --- not
just raw usage --- in the Prometheus text format:

\begin{shell}
docker compose -f docker-compose.grafana.yml up -d
# Grafana on http://localhost:3000, datasource and dashboard provisioned
\end{shell}

Scores are exported as \code{_sum}/\code{_count} pairs rather than
pre-averaged gauges, so PromQL keeps the arithmetic honest across any grouping:

\begin{shell}
sum by (emitter) (copilotscope_quality_score_sum)
  / sum by (emitter) (copilotscope_quality_score_count)
\end{shell}

Everything carries an \code{emitter} label, so a regression in one assistant is
visible instead of averaged away. Per-session series are opt-in and capped,
because session identifiers are unbounded and would hand Prometheus a new time
series for every conversation ever held.

One trap is worth naming. The provisioned alert rules use ratios of gauges and
never \code{rate()} over a counter. The collector recomputes its Prometheus
families from a capped in-memory set, so a family declared as a counter can
decrease when sessions age out, and \code{rate()} reads a decrease as a counter
reset --- a rule built that way fires on eviction rather than on anything real.

\subsection{Alerts and the weekly digest}

A dashboard that has to be visited gets abandoned. An output that triggers a
decision gets renewed. Configure \code{CopilotScope:Alerts} and the collector
compares each week against the one before it, by repository, assistant and
model, and posts a webhook when a cohort's mean composite falls.

What deliberately does \emph{not} fire is the interesting part: a drop that came
with a \emph{confidence} drop is reported as a changed measurement basis rather
than as a regression. The composite renormalizes over the components that have
data, so a cohort that stopped reporting a signal is being measured differently,
not performing worse. Sending a team to hunt a change that never happened is how
an alert channel gets muted, and a muted channel is worth less than no channel
because the team believes they have coverage.

\subsection{Turning the score from an opinion into a measurement}

This is the most valuable thing a reader of this manual can contribute.

\begin{shell}
# collector appsettings.json
"CopilotScope": { "Labelling": { "Enabled": true } }
\end{shell}

A \emph{Rate this session} panel appears on the session view. Rate real
sessions, not seeded ones. Then:

\begin{shell}
GET /api/labels/export > calibration/labels.json
POST /api/calibration/report   # on the judge agent, :5400
\end{shell}

The report endpoint is pure arithmetic: deterministic, free, and runnable in CI.
The research plan asks for fifty to a hundred labelled sessions before the
$\kappa$ means anything.

\subsection{Working from a clone}

\begin{shell}
git clone https://github.com/konradcinkusz/copilot-scope.git
cd copilot-scope
dotnet build
dotnet run --project src/CopilotScope.AppHost     # Aspire: everything
dotnet test                                        # no Docker needed
\end{shell}

The two opt-in agent services sit behind a compose profile, so a plain
\code{docker compose up} does not start them:

\begin{shell}
docker compose --profile agents up -d --build
\end{shell}

\subsection{When nothing shows up}

\begin{shell}
copilotscope doctor
\end{shell}

It checks Docker and the containers, the collector's health document, whether
the deployment is exposed without a key, what each assistant's settings file
actually says, whether an exported \codewrap{OTEL_EXPORTER_OTLP_ENDPOINT} is
overriding that file, and how many transcripts are sitting on disk unimported.

The classic causes, in the order they actually occur:

\begin{enumerate}[leftmargin=*]
    \item VS Code was not reloaded after the settings changed.
    \item Inline completions only --- no chat or agent turn happened.
    \item For Claude Code: the master switch or the logs exporter is missing.
    \item An exported environment variable beats the settings file for anything
    launched from that shell.
    \item Enterprise managed settings pin the endpoint to a corporate collector.
    Managed values always win.
    \item Cowork was pointed at the base endpoint rather than the full
    \code{/v1/logs} path, or the app was not restarted.
\end{enumerate}

% ============================================================================
\section{Operating it}
% ============================================================================

\subsection{What blocks a merge, and what a tag releases}

\begin{figure}[htbp]
\centering
\includegraphics[width=\linewidth]{../diagrams/rendered/d2-ci-gates.pdf}
\caption{For most of this project's life the images were built only on a release
tag. That is the gap the container job closes.}
\label{fig:ci}
\end{figure}

The container job earned its place in its first run: it found two pre-existing
bugs that had made every image unbuildable. The first was a shared project
missing from every build context, which the solution build could not see because
there the project sits exactly where the reference points. The second, reachable
only once the first was fixed, was two web projects resolving their
\code{appsettings.json} to the same publish path. Neither is exotic; both were
invisible because nothing built the images until a release was already being
cut.

\subsection{Retention, and what grows}

History grows until a policy is set. Retention is off by default:

\begin{shell}
"CopilotScope": { "History": {
  "RetentionDays": 0,        // 0 = keep everything; N = delete idle N+ days
  "BaselineDays": 30,        // window the percentile rank is computed over
  "BaselineMaxSamples": 5000
} }
\end{shell}

\code{BaselineDays} is what makes a percentile comparable over time: the rank is
against a fixed trailing window of user sessions, not against whatever happened
to still be in memory when you looked.

\subsection{Archiving the vendor's own numbers before they expire}

GitHub's Copilot Metrics API serves a 28-day rolling window and nothing older.
With \codewrap{CopilotScope:VendorMetrics} configured --- an organisation or
enterprise, a read-only token, and Postgres --- the collector polls daily and
keeps the window indefinitely, keyed per day so a re-poll overwrites rather than
accumulating duplicates.

The full response document is stored verbatim alongside the extracted counts,
because the vendor deletes the original: a parser that kept only what it
understood today would silently discard history that cannot be re-fetched.

The token needs exactly one read scope and nothing else. Nothing in this feature
writes, and a token with write scopes has more blast radius than the feature has
use for. Organisation and team level only: the API can be asked for a
per-developer breakdown, and this deliberately does not ask.

This is context, not the measurement. Counting usage still does not tell you
whether the tooling is helping. What the archive buys is a denominator --- the
sentence ``seats went up forty per cent in March'' is what makes a quality trend
readable --- and something useful on day one for an organisation with seats and
no instrumentation anywhere.

\subsection{Did the code actually ship?}

Every signal above stops at the session boundary. A session can run cleanly,
have its edits accepted, and still produce a change nobody merged --- which is
the same critique this project levels at usage counters, turned back on itself.

Outcome linkage closes that loop by joining sessions to pull requests. Point a
GitHub webhook at \code{POST /api/outcomes/github} with a secret; deliveries are
verified by HMAC over the raw body, and without a secret the endpoint is not
mapped at all, because it writes into the very data the score is about to be
validated against.

The join is a heuristic --- telemetry carries a repository and a branch, never a
commit --- so every link ships with its confidence and the reason for it, and a
repository-only match is shown but excluded from any correlation. Outcomes sit
\emph{beside} the score and are deliberately not folded into it, for the reason
given in Section~7.4.

% ============================================================================
\section{Where everything lives}
% ============================================================================

\begin{longtable}{@{}p{0.42\linewidth}p{0.54\linewidth}@{}}
\toprule
\textbf{Path} & \textbf{What is there} \\
\midrule
\endhead
\codepath{install.sh}, \code{install.ps1} & The one-command installer \\
\codepath{scripts/copilotscope} & The control script: up, connect, import, demo,
probe, doctor \\
\codepath{docker-compose.ghcr.yml} & The canonical pull-and-run stack \\
\codepath{docs/TUTORIAL.md} & The manual walkthrough, per assistant, plus
troubleshooting \\
\codepath{docs/SIGNAL_COVERAGE.md} & The full per-emitter signal matrix \\
\codepath{docs/PRIVACY.md} & GDPR Article 30 data map, retention, a template
works-agreement annex \\
\codepath{docs/CALIBRATION.md} & Method, scale anchors, how to read a bad result \\
\codepath{docs/DIAGRAMS.md} & Every diagram in this manual, rendered by GitHub \\
\codepath{docs/diagrams/} & The same diagrams as reusable sources \\
\codepath{docs/papers/} & This manual, both language editions \\
\codepath{docs/tutorials/} & Step-by-step tutorials, both languages \\
\codepath{docs/architecture/} & Architecture decision records and reviews \\
\codepath{SECURITY.md} & The trust model, scope by scope \\
\codepath{GOVERNANCE.md} & Who maintains it, what is stable, what happens if the
maintainer stops \\
\codepath{research/articles/} & The formal treatment of the scoring model, and ten
thesis topics derived from it \\
\bottomrule
\caption{The map. Where this manual and another document disagree, the other
document is the source of truth and this one is stale.}
\end{longtable}

% ============================================================================
\section{Conclusion}
% ============================================================================

The claim this project makes is narrow and worth restating precisely. Counting
how much an AI assistant was used does not tell you whether it helped. Scoring
the session does --- imperfectly, on evidence that varies by assistant, with a
formula that has not yet been validated against human judgement, and never about
a person.

Three commercial products now score AI coding-assistant sessions, which is the
strongest external evidence the thesis has: three companies with far better
market research concluded the problem is real. What remains distinctive here is
not the idea but the terms. Nothing leaves your infrastructure. The formula is
published and re-derivable by hand. The tool cannot rank developers, and that is
enforced by tests rather than promised in a document.

The honest next step is calibration, and it needs labelled sessions rather than
more code. The machinery is built, the export format is pinned by a test, and
the panel is one configuration flag away. Until then the composite is an opinion
with a confidence interval attached --- which is a useful thing, provided nobody
reads it as a measurement.

\section*{Reproducibility}

\begin{shell}
git clone https://github.com/konradcinkusz/copilot-scope.git
cd copilot-scope
npm ci                    # the pinned mermaid CLI
npm run render:diagrams   # docs/diagrams/*.mmd -> vector PDFs
pdflatex docs/papers/copilotscope-manual.tex   # twice, for the table of contents
\end{shell}

Both language editions are built together by
\code{.github/workflows/build-docs-pdf.yml}, on demand. PDFs are build output
and are never committed. \code{scripts/check-diagrams.mjs} fails the build if a
diagram disagrees with its embedded copy, and
\code{scripts/check-doc-parity.mjs} fails it if one language edition is edited
without the other.

\end{document}
